\begin{figure}[htbp]
\centering
\begin{tikzpicture}

% --- Top Row: Three Distribution Plots ---

% Prior Distribution
\begin{axis}[
    name=prior,
    width=5.5cm, height=5cm,
    xlabel={Parameter $\theta$},
    ylabel={Probability Density},
    xmin=-2, xmax=6,
    ymin=0, ymax=0.45,
    grid=major,
    grid style={gray!30},
    axis on top,
    title={\bfseries\color{UofSCAtlantic} Prior $p(\theta)$},
    title style={at={(0.5,1.0)}, anchor=south}
]

% Prior: Gaussian centered at 1 with wide variance
\addplot[fill=UofSCAtlantic!30, draw=UofSCAtlantic, very thick, domain=-2:6, samples=150]
    {1/(1.5*sqrt(2*pi)) * exp(-(x-1)^2/(2*1.5^2))} \closedcycle;

% Mean marker
\draw[UofSCAtlantic, thick, dashed] (1, 0) -- (1, 0.28);
\node[UofSCAtlantic, font=\footnotesize, above] at (1, 0.28) {$\mu_0=1$};

% Variance indicator
\draw[UofSCAtlantic, thick, |<->|] (1, 0.12) -- (2.5, 0.12);
\node[UofSCAtlantic, font=\footnotesize, above] at (1.75, 0.12) {$\sigma_0$};

\end{axis}

% Likelihood Function
\begin{axis}[
    name=likelihood,
    at={(prior.south east)},
    anchor=south west,
    xshift=0.8cm,
    width=5.5cm, height=5cm,
    xlabel={Parameter $\theta$},
    ylabel={},
    xmin=-2, xmax=6,
    ymin=0, ymax=0.45,
    grid=major,
    grid style={gray!30},
    axis on top,
    title={\bfseries\color{UofSCHorseshoe} Likelihood $p(x|\theta)$},
    title style={at={(0.5,1.0)}, anchor=south}
]

% Likelihood: Gaussian centered at 3.5 (where data suggests theta should be)
\addplot[fill=UofSCHorseshoe!30, draw=UofSCHorseshoe, very thick, domain=-2:6, samples=150]
    {1/(1.0*sqrt(2*pi)) * exp(-(x-3.5)^2/(2*1.0^2))} \closedcycle;

% Peak marker
\draw[UofSCHorseshoe, thick, dashed] (3.5, 0) -- (3.5, 0.42);
\node[UofSCHorseshoe, font=\footnotesize, above] at (3.5, 0.42) {$x_{obs}=3.5$};

% Annotation
\node[UofSCHorseshoe, font=\footnotesize, align=center] at (0.5, 0.32) {Data\\Evidence};

\end{axis}

% Posterior Distribution
\begin{axis}[
    name=posterior,
    at={(likelihood.south east)},
    anchor=south west,
    xshift=0.8cm,
    width=5.5cm, height=5cm,
    xlabel={Parameter $\theta$},
    ylabel={},
    xmin=-2, xmax=6,
    ymin=0, ymax=0.45,
    grid=major,
    grid style={gray!30},
    axis on top,
    title={\bfseries\color{UofSCGarnet} Posterior $p(\theta|x)$},
    title style={at={(0.5,1.0)}, anchor=south}
]

% Posterior: Narrower Gaussian between prior and likelihood means
% Using Bayesian update: posterior mean = (sigma_1^2 * mu_0 + sigma_0^2 * x) / (sigma_0^2 + sigma_1^2)
% posterior_mean = (1*1 + 2.25*3.5) / (2.25 + 1) = 8.875/3.25 = 2.73
% posterior_var = 1/(1/2.25 + 1/1) = 1/1.44 = 0.69, sigma = 0.83
\addplot[fill=UofSCGarnet!30, draw=UofSCGarnet, very thick, domain=-2:6, samples=150]
    {1/(0.83*sqrt(2*pi)) * exp(-(x-2.73)^2/(2*0.83^2))} \closedcycle;

% Posterior mean marker
\draw[UofSCGarnet, thick, dashed] (2.73, 0) -- (2.73, 0.49);
\node[UofSCGarnet, font=\footnotesize, above] at (2.73, 0.49) {$\mu_{post}$};

% Show it's between prior and likelihood
\draw[UofSCCongaree, thin, dotted] (1, 0.35) -- (1, 0.38);
\draw[UofSCCongaree, thin, dotted] (3.5, 0.35) -- (3.5, 0.38);
\draw[UofSCCongaree, thin, <->] (1, 0.36) -- (3.5, 0.36);
\node[UofSCCongaree, font=\tiny, above] at (2.25, 0.36) {Compromise};

\end{axis}

% --- Bottom Row: Process Flow Diagram ---
\begin{scope}[yshift=-6.5cm]

% Flow boxes
\node[draw=UofSCAtlantic, fill=UofSCAtlantic!15, minimum width=3cm, minimum height=1.2cm,
      align=center, font=\small, sharp corners] (priorbox) at (1.5, 0)
      {\textbf{Prior Belief}\\$p(\theta)$};

\node[draw=UofSCHorseshoe, fill=UofSCHorseshoe!15, minimum width=3cm, minimum height=1.2cm,
      align=center, font=\small, sharp corners] (databox) at (6.5, 0)
      {\textbf{New Data}\\$x_{obs}$};

\node[draw=UofSCGarnet, fill=UofSCGarnet!15, minimum width=3.5cm, minimum height=1.2cm,
      align=center, font=\small, sharp corners] (posteriorbox) at (12, 0)
      {\textbf{Updated Belief}\\$p(\theta|x)$};

% Multiplication node
\node[draw=UofSCBlack, fill=white, circle, minimum size=0.8cm, font=\large] (mult) at (4, 0) {$\times$};

% Proportionality node
\node[draw=UofSCBlack, fill=white, minimum width=1.5cm, minimum height=0.8cm,
      font=\small, sharp corners] (prop) at (9, 0) {$\propto$};

% Arrows
\draw[->, thick, UofSCBlack] (priorbox.east) -- (mult.west);
\draw[->, thick, UofSCBlack] (databox.west) -- (mult.east);
\draw[->, thick, UofSCBlack] (mult.east) ++(0.3,0) -- (prop.west);
\draw[->, thick, UofSCBlack] (prop.east) -- (posteriorbox.west);

% Likelihood label above data arrow
\node[UofSCHorseshoe, font=\footnotesize] at (5.25, 0.5) {Likelihood $p(x|\theta)$};

% Bayes' Theorem equation
\node[font=\normalsize, align=center] at (6.75, -1.5)
    {$\displaystyle p(\theta|x) = \frac{p(x|\theta) \cdot p(\theta)}{p(x)}$};
\node[font=\small, UofSC70Black] at (6.75, -2.2)
    {Posterior $\propto$ Likelihood $\times$ Prior};

% Iteration arrow (feedback)
\draw[->, thick, UofSCRose, dashed] (posteriorbox.south) -- ++(0, -0.7) -| (priorbox.south)
    node[pos=0.25, below, font=\footnotesize, UofSCRose] {Next iteration: posterior becomes new prior};

\end{scope}

\end{tikzpicture}
\caption{Bayesian updating process showing how prior beliefs are combined with data to form posterior estimates. \textbf{Top row}: The prior distribution (Atlantic blue, $\mu_0=1$, broad uncertainty) represents initial knowledge before observing data. The likelihood function (Horseshoe green) quantifies how probable the observed data $x_{obs}=3.5$ is for each possible parameter value. The posterior distribution (Garnet) is the normalized product, representing updated knowledge---narrower than the prior (reduced uncertainty) with mean shifted toward the data. \textbf{Bottom}: Schematic of Bayes' theorem. The dashed Rose arrow shows sequential updating: each posterior becomes the prior for the next observation. This framework underlies Kalman filtering, recursive estimation, and probabilistic inference in control systems.}
\label{fig:bayes_update}
\end{figure}
